% Nouveau EGA V (cote 25): archive cote 25, site page 87 (PDF page 88).
% typescript. Draft compared: none.
% AI-transcribed by Claude Opus 5.5 (claude-opus-5-5, adaptive thinking, effort high), Message Batches API, batch ega25b (msgbatch_01A7n5YuXubUPNK1zsepi8RG), 2026-11-13; not yet checked by eye.
\scanpage{25/87}
\mspage{81}

suivant qu'on suppose  dim X \hand{$\geqslant$} m , ou  dim X \hand{$\leqslant$} m , voire dim X \hand{$<$} m.
Par la suite, nous regardons la $C_{\hand{y}}\hand{\subset} P_{k(\hand{y})}$\hand{)} , correspondant au point
générique \hand{$y$} de  Grass$^{m+1}$ , et nous dispensons  de faire l'interprétation
des résultats obtenus en termes de  ``presque tous les points'' ...

Tout d'abord, nous avons déjà noté \var{\xout{\illegible{1}}}{}{typewriter strike-out}{ega25_25_087-v0001}dans 5.3. (et rattrapage du
cas général dans nº 12)  que $C_{\hand{y}}$  coupe ``régulièrement'' $X_{\hand{y}}$ , de façon
plus pré\var{\xout{\unsure{s}}}{}{typewriter strike-out}{ega25_25_087-v0002}cise et plus générale, pour tout F quasi-cohérent sur X,
la section $\phi_{\hand{y}}^{\hand{(m+1)}}$ du \var{\inserted[typed]{Module loc libre}}{Module loc libre}{insertion}{ega25_25_087-v0003}\var{\xout{fibré vectoriel}}{}{typewriter strike-out}{ega25_25_087-v0004} de rang  m+1  sur  $X_{k(\hand{y})}$  dont le
schéma des zéros est \var{\over{C}{$C$}}{$C$}{writing over}{ega25_25_087-v0005}$_{\hand{y}}$ ; est F-régulière. En vertu de 10.2. , cela
implique par exemple que le morphisme  $\widetilde{X}(\hand{C_y})\to X_{k(\hand{y})}$ identifie $\widetilde{X}(\hand{C_y})$ au préschéma
déduit de $X_{k(\hand{y})}$ par éclatement de $f^{-1}(C_{\hand{y}})= X\times_P C_{\hand{y}}$ ;  \var{\xout{\illegible{2}}}{}{typewriter strike-out}{ega25_25_087-v0006} dans le cas
où  dim f(X) \hand{$\leqslant$} m ,  on  aura même  $f^{-1}(C_{\hand{y}}) = \emptyset$ , et  par suite
$\widetilde{X}(C_{\hand{y}}) \xrightarrow{\sim} X_{k(\hand{y})}$  est un \uline{isomorphisme} , (et en fait, la \hand{\unsure{)}}projection conique
restreinte  est alors définie sur X tout entier à priori). Il se pose
ensuite la question de la dimension des fibres de $p_{C_{\hand{y}}}$ : $\widetilde{X}(C_{\hand{y}}) \longrightarrow Q(C_{\hand{y}})$,
et de la platitude de ce morphisme. On trouve:
\uline{Proposition} 14.1.  \var{\xout{a) Si dim X}}{}{typewriter strike-out}{ega25_25_087-v0007} Supposons X irréductible,\var{\over{\unsure{e}}{\unsure{ou}}}{\unsure{ou}}{writing over}{ega25_25_087-v0008}plus
généralement, que  pour toute composante irréductible $X_i$ de X, la fibre
de $X_i$ en le point $f(x_i)$  ($x_i$ = point générique de $X_i$) ait une dimension d
(indépendante de i), ce qui a lieu par exemple \var{\inserted{(avec d=0)}}{(avec d=0)}{insertion}{ega25_25_087-v0009} si f:X$\to$P est quasi-fini\hand{,} \var{\margin{\var{\deleted{avec d=0}}{}{deletion}{ega25_25_087-v0011}}}{\placedmargin{\var{\deleted{avec d=0}}{}{deletion}{ega25_25_087-v0011}}}{marginal note placed}{ega25_25_087-v0010}
Alors :     a)  Si  dim f(X) \hand{$>$} m, alors la dimension  des fibres de
$p_C$  : $\widetilde{X}(C_{\hand{y}}) \longrightarrow \var{\over{Q}{Q}}{Q}{writing over}{ega25_25_087-v0012}(C_{\hand{y}})$  sont toutes égales à  dim X - m .   b)  Si dim X
$\leqslant$ m, \var{\inserted{et si les fibres non vides des $X_i$ sur P sont \unsure{tts} de dimen d, alors}}{et si les fibres non vides des $X_i$ sur P sont \unsure{tts} de dimen d, alors}{insertion}{ega25_25_087-v0013} les fibres de $p_C$ \var{\inserted{non vides}}{non vides}{insertion}{ega25_25_087-v0014} \var{\inserted{$\colon X_{k(y)}\to Q(C_y)$}}{$\colon X_{k(y)}\to Q(C_y)$}{insertion}{ega25_25_087-v0015} sont toutes de dimension \var{\over{0}{d}}{d}{writing over}{ega25_25_087-v0016}, \var{\deleted{i.e. $p_C$ est quasi-%
fini}}{}{deletion}{ega25_25_087-v0017} (donc $p_C$ est fini \var{\inserted{\unsure{resp. quasi-fini}}}{\unsure{resp. quasi-fini}}{insertion}{ega25_25_087-v0018} \var{\deleted{si X est propre sur \unsure{k$(y)$} et}}{}{deletion}{ega25_25_087-v0019} \var{\inserted{$f\colon X\to P$ est fini \unsure{resp. quasi-fini}}}{$f\colon X\to P$ est fini \unsure{resp. quasi-fini}}{insertion}{ega25_25_087-v0020}].
\var{\margin{\unsure{resp. quasi-fini}}}{\placedmargin{\unsure{resp. quasi-fini}}}{marginal note placed}{ega25_25_087-v0021}

On \var{\deleted{se ramène aussitôt, en décomposant X en ses composantes irréd.,}}{}{deletion}{ega25_25_087-v0022}
\var{\deleted{au cas où X est irréductible}}{}{deletion}{ega25_25_087-v0023}. Dans le cas a), \var{\xout{si il y avait une fibre}}{}{typewriter strike-out}{ega25_25_087-v0024}
\var{\xout{correspondant à un}}{}{typewriter strike-out}{ega25_25_087-v0025} \var{\xout{xxxxxxxxxxxx}}{}{typewriter strike-out}{ega25_25_087-v0026} on sait déjà (je l'espère !) que pour \uline{tout}  point \hand{$y$} de

\var{\margin{87}}{\placedmargin{87}}{marginal note placed}{ega25_25_087-v0027}
